\section{Finite-molecule operation}\label{sec:finite}
Stable equilibria of the rate equations become usable labels only if a finite number of
molecules stays near them, returns after a disturbance, and can be told apart. This section
formulates such a contract for the full stochastic model of \eqref{eq:source} and proves
computable sufficient conditions. No diffusion approximation and no stochastic
quasi-steady-state reduction is used; such reductions need additional hypotheses
\cite{KimJosicBennett2014}.

\subsection{The Markov chain and exact lumping}
At system size \(\Omega>0\) the state is a vector \(N\) of molecule counts, and
\(X=N/\Omega\) is the vector of concentrations. A binding reaction with constant \(k\) and
reactant counts \(N_i,N_j\) has propensity \(kN_iN_j/\Omega\); the two reactants are always
distinct species, so no combinatorial correction arises. A unimolecular reaction has
propensity \(kN_i\). With \(a_\rho(x)\) the mass-action rate of reaction \(\rho\) and
\(\nu_\rho\) its reaction vector, the generator is
\begin{equation}\label{eq:generator}
 \cL_\Omega\varphi(x)=\Omega\sum_\rho a_\rho(x)\,
 \bigl[\varphi(x+\nu_\rho/\Omega)-\varphi(x)\bigr]
\end{equation}
\cite{AndersonKurtz2015}. The totals \eqref{eq:totals} are conserved, so for \(\Omega\) such
that \(\Omega\ET,\Omega\FT,\Omega\ST\) are integers the chain lives on a finite set, and it is
nonexplosive. We use the class chart of the proof of Theorem~\ref{thm:upper}: the \(d\)
coordinates are the \(S_I\) with \(I\ne\varnothing\) and all complexes, and
\(S_\varnothing,E,F\) are recovered from the totals. In this chart \(\nu_\rho\in\R^d\), \(H\)
denotes the vector field, \(J\) its Jacobian, and \(\|\cdot\|\) the Euclidean norm.

\begin{proposition}[Exact stochastic reduction]\label{prop:lump}
For the symmetric lift \eqref{eq:lift}, and for every initial distribution, the aggregated
process \(\Pi N(t)\) is a Markov chain whose generator is that of the sequential chain at the
same \(\Omega\). Consequently any event defined by the aggregated path has the same
probability in both models.
\end{proposition}
\begin{proof}
The reactions of the cube that induce a given jump of \(\Pi N\) are those of one type at one
level. Their propensities add up to
\(\sum_{|\tl e|=k}\frac{a_{k+1}}{n-k}N_EN_{S_{\tl e}}/\Omega
=a_{k+1}N_E\bigl(\sum_{|I|=k}N_{S_I}\bigr)/\Omega\) for kinase binding, similarly for
phosphatase binding, and to a level constant times an aggregated complex count for the
unimolecular reactions. These sums depend on \(N\) only through \(\Pi N\), which is the
criterion for strong lumpability \cite{KemenySnell1960}.
\end{proof}
The reduction is exact, not asymptotic, but it is limited to observations that factor through
\(\Pi\) and to the lifted rates; independently perturbed site-specific rates or site-specific
readouts are not covered.

\subsection{Retention and recovery}
\begin{theorem}[Computable finite memory]\label{thm:finite}
Let \(x^*\) be a positive hyperbolic sink with chart Jacobian \(J\). Let \(P\) be a symmetric
matrix and \(\underline p,\overline p,\mu,K,\varrho>0\) constants such that
\begin{equation}\label{eq:lyap-hyp}
 \underline pI\le P\le\overline pI,\qquad -(J^TP+PJ)\ge\mu I,\qquad
 \|H(x^*+h)-Jh\|\le K\|h\|^2\ \ (\|h\|\le\varrho),
\end{equation}
the ball \(\|h\|\le\varrho\) lies in the positive class, and
\(\varrho\le\mu/(4\overline pK)\). Put
\[
 V(x)=h^TPh\ \ (h=x-x^*),\qquad v=\underline p\varrho^2,\qquad
 \lambda=\frac{\mu}{2\overline p},\qquad D=\{V<v\},
\]
and let \(C\ge\sup_{x\in D}\sum_\rho a_\rho(x)\,\nu_\rho^TP\nu_\rho\). Let \(\tau\) be the
first exit time of \(X\) from \(D\). Then for every \(\Omega\), every \(X_0\in D\) and every
\(t\ge0\),
\begin{align}
 \Pr(\tau\le t)&\le\frac{V(X_0)+Ct/\Omega}{v},\label{eq:retention}\\
 \mathbb E\bigl[V(X_t)\,\mathbf 1_{\{\tau>t\}}\bigr]
 &\le e^{-\lambda t}V(X_0)+\frac{C}{\lambda\Omega}.\label{eq:recovery}
\end{align}
In particular, with the preparation set \(\mathsf B=\{V\le v/4096\}\) and the disturbance set
\(\mathsf A=\{V\le v/1024\}\), for every \(X_0\in\mathsf A\) and \(t\ge8/\lambda\),
\begin{equation}\label{eq:recfailure}
 \Pr\bigl(\tau\le t\ \text{or}\ X_t\notin\mathsf B\bigr)
 \le\frac1{1024}+\frac{Ct}{\Omega v}+\frac1{256}+\frac{4096\,C}{\lambda\Omega v}.
\end{equation}
\end{theorem}
\begin{proof}
Since \(V\) is quadratic, \eqref{eq:generator} gives the exact identity
\begin{equation}\label{eq:LV}
 \cL_\Omega V(x)=2h^TPH(x)+\frac1\Omega\sum_\rho a_\rho(x)\,\nu_\rho^TP\nu_\rho ,
\end{equation}
with no remainder. On \(D\) we have \(\underline p\|h\|^2\le V<\underline p\varrho^2\), so
\(\|h\|<\varrho\), and by \eqref{eq:lyap-hyp}
\(2h^TPH\le-\mu\|h\|^2+2\overline pK\|h\|^3\le-\tfrac\mu2\|h\|^2\le-\lambda V\). Hence
\(\cL_\Omega V\le-\lambda V+C/\Omega\) on \(D\). Dynkin's formula for the stopped process gives
\(\mathbb EV(X_{t\wedge\tau})\le V(X_0)+Ct/\Omega\). At the exit time the chain has jumped out
of \(D\), so \(V(X_\tau)\ge v\) whatever the overshoot, and Markov's inequality gives
\eqref{eq:retention}. Applying Dynkin's formula to \(e^{\lambda s}V(X_s)\) gives
\[
 \mathbb E\bigl[e^{\lambda(t\wedge\tau)}V(X_{t\wedge\tau})\bigr]
 \le V(X_0)+\frac C\Omega\int_0^te^{\lambda s}\,\dd s ,
\]
and \(e^{\lambda t}V(X_t)\mathbf 1_{\{\tau>t\}}\) is bounded by the integrand on the left, which
gives \eqref{eq:recovery}. Finally, by Markov's inequality,
\(\Pr(\tau>t,\ V(X_t)>v/4096)\le(4096/v)\bigl(e^{-\lambda t}v/1024+C/(\lambda\Omega)\bigr)\), and
\(4e^{-8}<1/256\); adding \eqref{eq:retention} with \(V(X_0)\le v/1024\) gives
\eqref{eq:recfailure}.
\end{proof}
The scalar form of the identity \eqref{eq:LV} and the multiplicity identity used in
Proposition~\ref{prop:lump} are among the machine-checked algebraic facts
(Appendix~\ref{app:certificates}).

\begin{corollary}[Operational contract]\label{cor:contract}
Let \(x_1^*,\dots,x_{N_{\mathrm{lab}}}^*\) be hyperbolic sinks in one class with data
\((P_i,\varrho_i,v_i,\lambda_i,C_i)\) as in Theorem~\textup{\ref{thm:finite}}. Let
\(\ell\) be a linear readout and \(\varepsilon_{\mathrm{obs}}\ge0\) a bound for the observation
error, and assume that the intervals
\([\ell(x_i^*)-\|\ell\|\varrho_i-\varepsilon_{\mathrm{obs}},\
\ell(x_i^*)+\|\ell\|\varrho_i+\varepsilon_{\mathrm{obs}}]\) are pairwise disjoint. Put
\(\trec=\max_i8/\lambda_i\), \(T=100\,\trec\), and let \(\Omega\) be compatible with the totals
and satisfy
\begin{equation}\label{eq:budget}
 \Omega\ \ge\ \max_i\max\Bigl\{\frac{1000\,C_iT}{v_i},\ \frac{10^6C_i}{\lambda_iv_i},\
 \frac{1024\,d\,\overline p_i}{v_i},\ 1\Bigr\}.
\end{equation}
Then: \textup{(a)} each \(\mathsf B_i\) contains a lattice point of the class, obtained by
rounding the chart coordinates of \(x_i^*\); \textup{(b)} from every state of \(\mathsf B_i\)
the chain stays in \(D_i\) during \([0,T]\) with probability at least \(1-0.0013\), and on this
event the interval decoder returns \(i\) at all times; \textup{(c)} from every state of
\(\mathsf A_i\) the chain is in \(\mathsf B_i\) at time \(\trec\), without having left \(D_i\),
with probability at least \(1-0.009\); \textup{(d)} recovery from \(\mathsf A_i\) followed by
storage for time \(T\) fails with probability less than \(0.011\).
\end{corollary}
\begin{proof}
(a) Rounding each chart coordinate of \(\Omega x_i^*\) to a nearest integer changes \(x\) by at
most \(1/(2\Omega)\) per coordinate, so \(V_i\le\overline p_id/(4\Omega^2)\le v_i/4096\) by the
third term of \eqref{eq:budget} and \(\Omega\ge1\); the dependent counts are then integers by
compatibility and positive because the ball lies in the positive class. When \(x_i^*\) is
irrational one rounds a sufficiently accurate enclosure instead. (b) and (c) follow from
\eqref{eq:retention} and \eqref{eq:recfailure}: the terms are at most
\(1/4096+10^{-3}\) and \(1/1024+10^{-5}+1/256+0.0041\). On \(\{\tau_i>T\}\) the readout differs
from \(\ell(x_i^*)\) by at most \(\|\ell\|\varrho_i\). (d) is the Markov property at time
\(\trec\) and a union bound.
\end{proof}

\begin{remark}[Three limits of the contract]\label{rem:contract}
(a) Retention is continuous: an excursion out of \(D_i\) counts as a failure even if the chain
returns. (b) Recovery is from the local disturbance sets \(\mathsf A_i\), not from arbitrary
states. (c) Labels are \emph{prepared}. No protocol that switches the system from one label to
another with a common finite control is analysed, and the corollary does not describe a
complete write--store--read device. The individual guarantees (b) and (c) exceed \(0.99\); the
composite guarantee (d) is what may be quoted for recovery followed by storage.
\end{remark}

\begin{proposition}[A common clock does not change the budget]\label{prop:clock}
Multiply all rate constants by \(\varsigma>0\) and keep the totals, the matrices \(P_i\) and
the radii \(\varrho_i\). Then the hypotheses \eqref{eq:lyap-hyp} hold with
\((\mu_i,K_i,C_i,\lambda_i)\) replaced by \(\varsigma(\mu_i,K_i,C_i,\lambda_i)\), the times
\(\trec\) and \(T\) are divided by \(\varsigma\), and the right-hand side of \eqref{eq:budget}
is unchanged.
\end{proposition}
\begin{proof}
The generator becomes \(\varsigma\cL_\Omega\), so the new process is the old one run at time
\(\varsigma t\). The ratio \(\mu_i/K_i\), and hence the admissible radius and \(v_i\), are
unchanged, and \(C_iT\) and \(C_i/\lambda_i\) are invariant.
\end{proof}
Thus slowing a design to respect physical upper limits on rate constants changes the reported
times but not the molecule budget of the matched protocol \(T=100\,\trec\); at a fixed absolute
storage time the bound would change, and recovery would have to be charged on the same clock.

\section{Crowding, restoration, and observation}\label{sec:crowding}

\subsection{The exit barrier of crowded states}
The chain attains its \(n\) sinks in \cite{Companion} by splitting one equilibrium of
multiplicity \(2n-1\). For \(n\ge2\) and real numbers \(\xi_0<\dots<\xi_{2n-2}\), that paper
constructs rate constants \(\kappa(\delta)\), analytic in a small parameter \(\delta\ge0\),
with fixed totals, such that for \(\delta>0\) the class contains exactly the \(2n-1\)
equilibria \(X_j(\delta)\), with \(E=1+\delta\xi_j/2\) at \(X_j(\delta)\), all converging to
one point \(X_\circ\) as \(\delta\to0\); the \(X_j\) with \(j\) even are hyperbolic sinks, and
the Jacobian at \(X_j(\delta)\) has \(3n-1\) eigenvalues with real parts below a negative
constant and one eigenvalue
\begin{equation}\label{eq:split}
 \lambda_j(\delta)=-a\,\delta^{2n-2}\prod_{k\ne j}(\xi_j-\xi_k)\,(1+O(\delta)),\qquad a>0
\end{equation}
(Appendix~\ref{app:sequential}). We call this the \emph{crowded family}. By
Propositions~\ref{prop:lift} and~\ref{prop:lump} everything below applies verbatim to the
symmetric lift of the family to the cube and to events defined through the aggregation
\(\Pi\).

For a path \(\phi:[0,t]\to\) class chart, the reaction-path action is
\(\int_0^t\cI(\phi,\dot\phi)\,\dd s\), where \(\cI(x,\cdot)\) is the Legendre transform of the
Hamiltonian
\begin{equation}\label{eq:hamiltonian}
 \mathcal H(x,\theta)=\sum_\rho a_\rho(x)\bigl(e^{\theta\cdot\nu_\rho}-1\bigr),
\end{equation}
the large-deviation rate function of the chain \eqref{eq:generator}
\cite{ShwartzWeiss1995,FreidlinWentzell,AgazziDemboEckmann2018}.

\begin{theorem}[Local action and iterated exit law]\label{thm:action}
Consider the crowded family for some \(n\ge2\), and let \(i\) be even. For every
\(k\ge1\) and all sufficiently small \(\delta\) there are \(C^k\) coordinates \((z,w)\in
\R\times\R^{3n-1}\) on a fixed neighbourhood of \(X_\circ\) in the class, depending \(C^k\) on
\(\delta\), such that the rate equations read
\begin{equation}\label{eq:normalform}
 \dot z=f(z,\delta)=-a\prod_{j=0}^{2n-2}(z-\delta\xi_j)\,\bigl(1+O(|z|+\delta)\bigr),\qquad
 \dot w=\mathsf B(z,w,\delta)\,w ,
\end{equation}
with \(a\) as in \eqref{eq:split} and \(\mathsf B(0,0,0)\) Hurwitz; \(X_j(\delta)\) is the point
\((\delta\xi_j,0)\). Let \(l<\xi_i<r_{\mathrm c}\) be cuts with
\(\xi_{i-1}<l\) and \(r_{\mathrm c}<\xi_{i+1}\) (when these neighbours exist), let \(P_t\) be a
positive definite matrix with \(\mathsf B(0,0,0)^TP_t+P_t\mathsf B(0,0,0)<0\), let
\(1<\alpha_t<n\), and put
\[
 D_\delta=\bigl\{\delta l<z<\delta r_{\mathrm c},\ \ w^TP_tw<\delta^{2\alpha_t}\bigr\}.
\]
\textup{(i)} The infimum \(I_\delta\) of the action over all paths from \(X_i(\delta)\) to the
complement of \(D_\delta\) satisfies
\begin{equation}\label{eq:action}
 I_\delta=c_i\,\delta^{2n}\,(1+O(\delta)),\qquad
 c_i=\frac{2a}{b_0}\ \min_{s\in\{l,r_{\mathrm c}\}}\ \int_{\xi_i}^{s}\prod_{j}(y-\xi_j)\,\dd y\ >0,
\end{equation}
where \(b_0=\omega^TQ(X_\circ)\,\omega\), \(Q(x)=\sum_\rho a_\rho(x)\nu_\rho\nu_\rho^T\), and
\(\omega=\dd z\) at \(X_\circ\) is the left null covector of the Jacobian normalized by the
choice of \(z\). \textup{(ii)} Let \(\tau_\delta\) be the exit time from \(D_\delta\) of the
chain started at lattice points \(X_0^\Omega\to X_i(\delta)\). Then for every \(\zeta>0\),
\begin{equation}\label{eq:iterated}
 \lim_{\delta\downarrow0}\ \liminf_{\Omega\to\infty}\
 \Pr\Bigl(\Bigl|\frac{\log\tau_\delta}{\Omega\,\delta^{2n}}-c_i\Bigr|<\zeta\Bigr)=1 .
\end{equation}
\end{theorem}
The proof is in Appendix~\ref{app:action}. The inner limit is taken at fixed \(\delta\); the
theorem is not uniform along joint sequences \((\delta,\Omega)\), and it concerns the declared
tube \(D_\delta\), not the basin of attraction. The action in (i) is the full reaction-path
action: paths may vary all enzyme, substrate and complex coordinates, and the transverse bound
\(\alpha_t<n\) is what excludes a cheaper escape through those coordinates.

\begin{corollary}[Barrier, recovery and resolution]\label{cor:tradeoff}
In the crowded family, the linear recovery time \(\trec=1/|\lambda_i|\), the gap
\(\Delta R\) between neighbouring values of any observable with nonzero derivative along the
equilibrium curve, and the barrier satisfy
\[
 \trec=\Theta(\delta^{-(2n-2)}),\quad \Delta R=\Theta(\delta),\quad
 I_\delta=\Theta(\delta^{2n})
 \quad\Longrightarrow\quad
 I_\delta=\Theta\bigl(\trec^{-n/(n-1)}\bigr)=\Theta\bigl((\Delta R)^{2n}\bigr),
\]
with constants depending on the family.
\end{corollary}
\begin{proof}
The first relation is \eqref{eq:split}, the second is \cite[Corollary~8.4]{Companion} or
follows from the analyticity of \(X_j(\delta)\), and the third is \eqref{eq:action}.
\end{proof}
The linear-noise estimate \(\Var(z)\simeq b_0/(2\Omega|\lambda_i|)\) gives a squared
separation-to-noise ratio of order \(\Omega\delta^{2n}\), in agreement with the exponent in
\eqref{eq:action}; it is a consistency check, not a substitute for the theorem. These relations
describe \emph{one} mechanism by which many states can be expensive, namely crowding. They do
not say that every multi-label source is crowded: the four labels of
Section~\ref{sec:example} are well separated, and their cost has other origins.

\subsection{Observation}
\begin{proposition}[Bounded-error packing]\label{prop:packing}
Suppose label \(i\) has ideal observation \(y_i\in\prod_{j=1}^{d_{\mathrm{obs}}}[\alpha_j,\alpha_j+R_j]\), the
observed vector is \(y_i+e\) with an arbitrary error satisfying
\(|e_j|\le\varepsilon_j\), \(\varepsilon_j>0\), and a decoder identifies the label for every
admissible error. Then the number of labels satisfies
\[
 N_{\mathrm{lab}}\ \le\ \prod_{j=1}^{d_{\mathrm{obs}}}\Bigl(1+\Bigl\lfloor\frac{R_j}{2\varepsilon_j}\Bigr\rfloor\Bigr).
\]
\end{proposition}
\begin{proof}
The closed error boxes around different ideal observations must be disjoint. Let
\(m_j=1+\lfloor R_j/2\varepsilon_j\rfloor>R_j/2\varepsilon_j\) and cut the \(j\)-th range into
\(m_j\) consecutive intervals of length \(R_j/m_j<2\varepsilon_j\), assigning endpoints to one
side. Two ideal observations in the same product cell differ by less than \(2\varepsilon_j\)
in every coordinate, so their boxes share the midpoint. Hence each cell contains at most one
ideal observation.
\end{proof}
Variation of the readout within the safe sets \(D_i\) enlarges the boxes. The bound is of order
\((R/\varepsilon)^{d_{\mathrm{obs}}}\), which is why several independent channels can help; it is an
upper bound, not an attainment result, and correlated channels, or phosphoform fractions
constrained to a simplex, give less.

\begin{proposition}[Passive and bounded-gain observation]\label{prop:passive}
\textup{(a)} If an observation is a function of the chemical path, possibly randomized, that
does not feed back on the chemistry, then the law of every chemical exit time, and hence every
barrier of Theorem~\textup{\ref{thm:action}}, is the same for all such observations.
\textup{(b)} If equilibria satisfy \(\|x_i(\delta)-x_j(\delta)\|=O(\delta)\) and the
observation map has a Lipschitz constant bounded independently of \(\delta\), then the
observations differ by \(O(\delta)\), for any number of channels.
\end{proposition}
\begin{proof}
(a) The exit time is a functional of the chemical path, whose law is unchanged. (b) is the
Lipschitz inequality.
\end{proof}
A richer decoder can change which excursions count as failures; that defines a new event, which
needs its own retention proof. A reporter that binds the substrate is not passive: it
sequesters substrate and belongs in the reaction network. Thus multichannel readout is a
promising direction for well-spread states, but it does not repair crowding.

\section{Finite thermodynamic driving and robustness}\label{sec:driving}
The catalytic steps of \eqref{eq:source} are irreversible, which corresponds to infinite
driving by ATP hydrolysis. We show that the sinks survive in a thermodynamically consistent
model with finite driving, and that effective affinities present no obstruction; see
\cite{Hill1989,Qian2007,RaoEsposito2016} for background and \cite{LiLiaoOuyangLi2024} for a
study of finite driving in other biochemical switches.

\begin{theorem}[Finite-driving persistence]\label{thm:driving}
Let ATP, ADP and inorganic phosphate be chemostatted with activities \(1\), \(\vartheta\) and
\(\vartheta\), and replace \eqref{eq:source} on every edge by
\begin{equation}\label{eq:reversible}
 S_{\tl e}+E+\mathrm{ATP}\ \underset{b_e}{\overset{a_e}{\rightleftharpoons}}\ C_e
 \ \underset{d_e}{\overset{c_e}{\rightleftharpoons}}\ S_{\hd e}+E+\mathrm{ADP},
 \qquad
 S_{\hd e}+F\ \underset{\beta_e}{\overset{\alpha_e}{\rightleftharpoons}}\ Y_e
 \ \underset{h_e}{\overset{\gamma_e}{\rightleftharpoons}}\ S_{\tl e}+F+\mathrm P_{\mathrm i},
\end{equation}
with \(d_e=a_ec_e/b_e\) and \(h_e=\alpha_e\gamma_e/\beta_e\). \textup{(i)} There are standard
chemical potentials for which every reaction of \eqref{eq:reversible} satisfies local
detailed balance; hence all Wegscheider cycle conditions hold, including those for cycles
through different paths of the cube, and the chemical potential difference of ATP hydrolysis,
\(-2\log\vartheta\) in units of \(\kB T\), is the same for every turnover on every edge.
\textup{(ii)} The vector field of \eqref{eq:reversible} is \(H+\vartheta G\) with a polynomial
\(G\), and the totals \eqref{eq:totals} are still conserved. Consequently every finite set of
hyperbolic sinks of \eqref{eq:source} in one class continues, for all sufficiently small
\(\vartheta>0\), to hyperbolic sinks of \eqref{eq:reversible} in the same class, and disjoint
readout neighbourhoods remain disjoint.
\end{theorem}
\begin{proof}
(i) Set the standard potentials, in units of \(\kB T\), to zero for all free species
\(S_I,E,F\) and for the three chemostatted species, to \(\log(b_e/a_e)\) for \(C_e\) and to
\(\log(\beta_e/\alpha_e)\) for \(Y_e\). Then each of the four reaction pairs has ratio of
forward to backward constant equal to \(\exp(-\Delta\mu^\circ)\): for the binding steps by
definition, and for the catalytic steps because \(c_e/d_e=b_e/a_e\) and
\(\gamma_e/h_e=\beta_e/\alpha_e\). Around any closed cycle of reactions the increments of a
potential telescope, so the product of the ratios is one \cite{SchusterSchuster1989}. With the
stated activities the hydrolysis potential is \(\log[1/(\vartheta\cdot\vartheta)]\). (ii) The
forward propensities are unchanged, and the new reverse propensities are
\(\vartheta d_eES_{\hd e}\) and \(\vartheta h_eFS_{\tl e}\). The \(C^1\) convergence on compact
sets is immediate, and hyperbolic equilibria and their spectra persist by the implicit
function theorem.
\end{proof}

\begin{remark}
(a) The three notions that occur in this paper are different: effective square balance
\eqref{eq:balance-intro}, which concerns the irreversible network; detailed balance of the
closed reversible network, which holds for \eqref{eq:reversible} by construction; and the
maintained hydrolysis potential. An arbitrary effective affinity is compatible with (i)
because it compares kinetic pathways, not equilibrium constants. (b) The choice of zero
standard potentials for the phosphoforms is a normalization of convenience; other potentials
change \(d_e,h_e\) by bounded factors and leave (ii) unchanged. (c) The theorem concerns a
chemostatted system. A finite closed fuel pool needs a depletion argument, and a quantitative
design needs feasible chemostat concentrations and ATP throughput. (d) The finite-molecule
constants of Theorem~\ref{thm:finite} do not transfer automatically: the added reactions change
the drift, the noise and the equilibria, and the contract must be recomputed.
\end{remark}

The following estimate makes ``sufficiently small'' explicit, and also gives a robustness
margin for the rate constants.

\begin{proposition}[Explicit perturbation bound]\label{prop:robust}
Assume the hypotheses of Theorem~\textup{\ref{thm:finite}} and let
\(D'=\{V\le v/4\}\). Let \(G\) be a \(C^1\) field on \(D'\), tangent to the class, with
\(\|G\|\le b_G\) and \(\|DG\|\le d_G\) there, and let
\(\ell_0=\min\{v/(4\overline p),1\}\). If
\begin{equation}\label{eq:robust}
 0\le\vartheta\le\min\Bigl\{\frac{\mu}{8\overline p\,d_G},\ \frac{\mu\,\ell_0}{8\overline p\,b_G}\Bigr\},
\end{equation}
then \(H+\vartheta G\) has exactly one equilibrium in \(D'\), and it is a hyperbolic sink.
\end{proposition}
\begin{proof}
On \(D'\) we have \(\|h\|\le\varrho/2\). For the mass-action field the remainder is a
quadratic form bounded by \(K\|h\|^2\) through a bound on its bilinear form, so
\(\|DH(x^*+h)-J\|\le2K\|h\|\le K\varrho\), and
\(DH^TP+PDH\le-\mu I+2\overline pK\varrho I\le-\tfrac\mu2I\) on \(D'\). Integrating along the
segment from \(x^*\) gives \(h^TPH(x^*+h)\le-\tfrac\mu4\|h\|^2\). On \(\partial D'\),
\(\|h\|^2\ge v/(4\overline p)\), so \(\|h\|\ge\ell_0\) because \(\min\{t,1\}\le\sqrt t\), and
\(\vartheta|h^TPG|\le\vartheta\overline pb_G\|h\|\le\tfrac\mu8\|h\|^2\). Thus the perturbed
field points strictly inward on \(\partial D'\), and a zero exists by degree. Moreover
\(D(H+\vartheta G)^TP+PD(H+\vartheta G)\le-\tfrac\mu2I+2\vartheta\overline pd_GI\le-\tfrac\mu4I\)
on the convex set \(D'\), so the flow contracts the metric \(P\) there: the zero is unique and
its Jacobian is Hurwitz by Lyapunov's theorem.
\end{proof}
For independent relative errors \(|\Delta k_\rho|\le\vartheta k_\rho\) of the rate constants one
may take \(b_G=\sum_\rho\sup a_\rho\|\nu_\rho\|\) and the analogous bound for \(d_G\); for the
reversible extension, \(G\) consists of the reverse catalytic reactions. In both cases rational
upper bounds on \(D'\) are computable, and the perturbed equilibrium stays in \(D'\), hence
within the decoder interval of Corollary~\ref{cor:contract}.
